\documentclass[12pt,a4paper,oneside]{article}
\usepackage[brazil]{babel}
\usepackage[utf8]{inputenc}
\usepackage{olymp}
\usepackage{color}
\usepackage[dvipsnames]{xcolor}
\usepackage{listings}

\setcounter{problem}{0}

\begin{document}

\begin{problem}{Fatorial por extenso}{fatex}{2}
 
O fatorial $N!$ de um número natural $N > 0$ é definido como a multiplicação de todos os naturais entre 1 e $N$, ou seja, $N!=N \times (N-1) \times (n-2) \times \dots \times 1$. Também pode ser definido recursivamente por $N! = N \times (N-1)!$, como na função do programa abaixo. Em todo caso, $0! = 1$, por definição.

\lstset{language=C++,
                basicstyle=\ttfamily,
                keywordstyle=\color{Mulberry}\ttfamily,
                stringstyle=\color{Maroon}\ttfamily,
                commentstyle=\color{YellowGreen}\ttfamily,
                morecomment=[l][\color{magenta}]{\#},
                basewidth = {.48em},
                showstringspaces=false
}
\begin{lstlisting}
//Retorna o valor de x!
long long int fatorial (int n) {
    if (n<=1)
        return 1;
    else
        return n * fatorial(n-1);
}

int main() {
    int n;
    
    cin >> n;
    cout << n << "! = ";
    cout << fatorial(n) << endl;
    return 0;
}

\end{lstlisting}

Para a entrada $5$, o programa imprime \texttt{5!} \texttt{= 120}.

Modifique a função fatorial do programa acima para escrever detalhadamente o cálculo do fatorial, como nos exemplos a seguir.

\Notes

Nesta questão, seu programa deve usar a função \texttt{main} dada acima exatamente como está, não a modifique! Sua tarefa é modificar a função \texttt{fatorial}.


\InputFile

A entrada contém um único número natural $N$. Restrições: $2 \leq N \leq 20$.

Note que seu programa não será testado com $N=0$ nem $N=1$.

\OutputFile

Seu programa deve gerar apenas uma linha de saída, contendo o cálculo detalhado do fatorial de $N$, como mostrado nos exemplos a seguir.

\Example

\setlength\exmpwidinf{6cm}
\setlength\exmpwidouf{10cm}
   

\begin{example}
\exmp{
5
}{
5! = 5 x 4 x 3 x 2 x 1 = 120
}%
\end{example}

\begin{example}
\exmp{
7
}{
7! = 7 x 6 x 5 x 4 x 3 x 2 x 1 = 5040
}%
\end{example}

%Comentários sobre o exemplo, se necessário.

\end{problem}

\end{document}
